% Options for packages loaded elsewhere
% Options for packages loaded elsewhere
\PassOptionsToPackage{unicode}{hyperref}
\PassOptionsToPackage{hyphens}{url}
%
\documentclass[
  ignorenonframetext,
  aspectratio=169,
]{beamer}
\newif\ifbibliography
\usepackage{pgfpages}
\setbeamertemplate{caption}[numbered]
\setbeamertemplate{caption label separator}{: }
\setbeamercolor{caption name}{fg=normal text.fg}
\beamertemplatenavigationsymbolsempty
% remove section numbering
\setbeamertemplate{part page}{
  \centering
  \begin{beamercolorbox}[sep=16pt,center]{part title}
    \usebeamerfont{part title}\insertpart\par
  \end{beamercolorbox}
}
\setbeamertemplate{section page}{
  \centering
  \begin{beamercolorbox}[sep=12pt,center]{section title}
    \usebeamerfont{section title}\insertsection\par
  \end{beamercolorbox}
}
\setbeamertemplate{subsection page}{
  \centering
  \begin{beamercolorbox}[sep=8pt,center]{subsection title}
    \usebeamerfont{subsection title}\insertsubsection\par
  \end{beamercolorbox}
}
% Prevent slide breaks in the middle of a paragraph
\widowpenalties 1 10000
\raggedbottom
\AtBeginPart{
  \frame{\partpage}
}
\AtBeginSection{
  \ifbibliography
  \else
    \frame{\sectionpage}
  \fi
}
\AtBeginSubsection{
  \frame{\subsectionpage}
}
\usepackage{iftex}
\ifPDFTeX
  \usepackage[T1]{fontenc}
  \usepackage[utf8]{inputenc}
  \usepackage{textcomp} % provide euro and other symbols
\else % if luatex or xetex
  \usepackage{unicode-math} % this also loads fontspec
  \defaultfontfeatures{Scale=MatchLowercase}
  \defaultfontfeatures[\rmfamily]{Ligatures=TeX,Scale=1}
\fi
\usepackage{lmodern}

\ifPDFTeX\else
  % xetex/luatex font selection
\fi
% Use upquote if available, for straight quotes in verbatim environments
\IfFileExists{upquote.sty}{\usepackage{upquote}}{}
\IfFileExists{microtype.sty}{% use microtype if available
  \usepackage[]{microtype}
  \UseMicrotypeSet[protrusion]{basicmath} % disable protrusion for tt fonts
}{}
\makeatletter
\@ifundefined{KOMAClassName}{% if non-KOMA class
  \IfFileExists{parskip.sty}{%
    \usepackage{parskip}
  }{% else
    \setlength{\parindent}{0pt}
    \setlength{\parskip}{6pt plus 2pt minus 1pt}}
}{% if KOMA class
  \KOMAoptions{parskip=half}}
\makeatother


\usepackage{longtable,booktabs,array}
\usepackage{calc} % for calculating minipage widths
\usepackage{caption}
% Make caption package work with longtable
\makeatletter
\def\fnum@table{\tablename~\thetable}
\makeatother
\usepackage{graphicx}
\makeatletter
\newsavebox\pandoc@box
\newcommand*\pandocbounded[1]{% scales image to fit in text height/width
  \sbox\pandoc@box{#1}%
  \Gscale@div\@tempa{\textheight}{\dimexpr\ht\pandoc@box+\dp\pandoc@box\relax}%
  \Gscale@div\@tempb{\linewidth}{\wd\pandoc@box}%
  \ifdim\@tempb\p@<\@tempa\p@\let\@tempa\@tempb\fi% select the smaller of both
  \ifdim\@tempa\p@<\p@\scalebox{\@tempa}{\usebox\pandoc@box}%
  \else\usebox{\pandoc@box}%
  \fi%
}
% Set default figure placement to htbp
\def\fps@figure{htbp}
\makeatother





\setlength{\emergencystretch}{3em} % prevent overfull lines

\providecommand{\tightlist}{%
  \setlength{\itemsep}{0pt}\setlength{\parskip}{0pt}}



 


\usepackage{booktabs}
\usepackage{adjustbox}
\usepackage{graphicx}
\usepackage{amsmath}
\usepackage{array}
\definecolor{DeepTeal}{HTML}{20393D}
\setbeamercolor{frametitle}{fg=white,bg=DeepTeal}
\setbeamerfont{frametitle}{series=\bfseries,size=\Large}
\setbeamertemplate{frametitle}{\nointerlineskip\begin{beamercolorbox}[wd=\paperwidth,ht=0.86cm,dp=0.23cm,leftskip=0.35cm]{frametitle}\insertframetitle\end{beamercolorbox}}
\setbeamertemplate{navigation symbols}{}
\setbeamertemplate{footline}[frame number]
\setbeamersize{text margin left=0.75cm,text margin right=0.75cm}
\setbeamerfont{itemize/enumerate body}{size=\large}
\setbeamerfont{itemize/enumerate subbody}{size=\normalsize}
\makeatletter
\@ifpackageloaded{caption}{}{\usepackage{caption}}
\AtBeginDocument{%
\ifdefined\contentsname
  \renewcommand*\contentsname{Table of contents}
\else
  \newcommand\contentsname{Table of contents}
\fi
\ifdefined\listfigurename
  \renewcommand*\listfigurename{List of Figures}
\else
  \newcommand\listfigurename{List of Figures}
\fi
\ifdefined\listtablename
  \renewcommand*\listtablename{List of Tables}
\else
  \newcommand\listtablename{List of Tables}
\fi
\ifdefined\figurename
  \renewcommand*\figurename{Figure}
\else
  \newcommand\figurename{Figure}
\fi
\ifdefined\tablename
  \renewcommand*\tablename{Table}
\else
  \newcommand\tablename{Table}
\fi
}
\@ifpackageloaded{float}{}{\usepackage{float}}
\floatstyle{ruled}
\@ifundefined{c@chapter}{\newfloat{codelisting}{h}{lop}}{\newfloat{codelisting}{h}{lop}[chapter]}
\floatname{codelisting}{Listing}
\newcommand*\listoflistings{\listof{codelisting}{List of Listings}}
\makeatother
\makeatletter
\makeatother
\makeatletter
\@ifpackageloaded{caption}{}{\usepackage{caption}}
\@ifpackageloaded{subcaption}{}{\usepackage{subcaption}}
\makeatother

\usepackage{bookmark}
\IfFileExists{xurl.sty}{\usepackage{xurl}}{} % add URL line breaks if available
\urlstyle{same}
\hypersetup{
  pdftitle={Exposição à Inteligência Artificial e Mercado de Trabalho no Brasil},
  pdfauthor={Manoel Brasil},
  hidelinks,
  pdfcreator={LaTeX via pandoc}}


\title{Exposição à Inteligência Artificial e Mercado de Trabalho no
Brasil}
\subtitle{Amostra, estatísticas descritivas e estratégia empírica}
\author{Manoel Brasil}
\date{}
\institute{Orientadora: Profa. Danyelle Karine Santos Branco\newline 14
de Maio}

\begin{document}
\frame{\titlepage}


\begin{frame}{Resumo da Apresentação}
\phantomsection\label{resumo-da-apresentauxe7uxe3o}
\begin{itemize}
\tightlist
\item
  \textbf{Etapa 1:} PNAD Contínua + índice ILO --- quem está exposto à
  IA no Brasil?
\item
  \textbf{Etapa 2:} Novo CAGED + DiD --- ocupações expostas mudaram após
  o ChatGPT?
\item
  \textbf{Etapa 3:} CAGED municipal + Anatel --- conectividade amplifica
  o efeito?
\item
  Para cada etapa: \textbf{amostra}, \textbf{estatísticas descritivas}
  e, quando houver inferência causal, \textbf{modelo econométrico}.
\end{itemize}
\end{frame}

\begin{frame}{Pergunta de Pesquisa e Desenho Geral}
\phantomsection\label{pergunta-de-pesquisa-e-desenho-geral}
\textbf{Pergunta:} a IA generativa já deixou sinais mensuráveis no
mercado de trabalho brasileiro?

O desenho empírico avança em três camadas:

\begin{itemize}
\tightlist
\item
  Primeiro, mapeio a exposição ocupacional à IA na força de trabalho
  brasileira.
\item
  Depois, testo se ocupações mais expostas tiveram trajetória diferente
  após o lançamento do ChatGPT.
\item
  Por fim, verifico se o efeito é maior onde a infraestrutura digital
  facilita a difusão da tecnologia.
\end{itemize}
\end{frame}

\section{Etapa 1: PNADc + ILO}\label{etapa-1-pnadc-ilo}

\begin{frame}{Etapa 1 --- Construção da Base}
\phantomsection\label{etapa-1-construuxe7uxe3o-da-base}
\begin{center}
\resizebox{0.78\textwidth}{!}{\scriptsize
\begin{tabular}{lr}
\toprule
Item & Valor \\
\midrule
Fonte & PNADc 3T/2025 (IBGE) + ILO WP140 \\
Observações na amostra & 207.901 \\
População representada & 97,8 milhões \\
Unidades federativas & 27 \\
Ocupações COD & 428 \\
Cobertura de score & 99,2\% \\
Match a 4 dígitos & 98,1\% \\
Alta exposição (G3-G4) & 9,8 milhões \\
\bottomrule
\end{tabular}
}
\end{center}

\vspace{0.2cm}

\begin{itemize}
\tightlist
\item
  A unidade de observação é o trabalhador ocupado na PNADc.
\item
  O peso amostral transforma a amostra em população representada.
\item
  A variável central é o score ocupacional do ILO Global Index,
  associado à ocupação COD por crosswalk hierárquico.
\item
  Tratamento de extremos: rendimento habitual limitado aos percentis
  ponderados 1 e 99.
\end{itemize}
\end{frame}

\begin{frame}{Etapa 1 --- Estatísticas Descritivas e Amostra}
\phantomsection\label{etapa-1-estatuxedsticas-descritivas-e-amostra}
\begin{center}
\resizebox{0.96\textwidth}{!}{\scriptsize
\begin{tabular}{lrrrrrr}
\toprule
Variável & Unidade & N & Média & Desv. Pad. & Min & Max \\
\midrule
Score de exposição ILO & 0-1 & 206.230 & 0,28 & 0,15 & 0,09 & 0,70 \\
Alta exposição (G3-G4) & \% & 207.901 & 10,00 & 30,00 & 0,00 & 100,00 \\
Idade & anos & 207.901 & 39,26 & 12,11 & 18,00 & 65,00 \\
Nível de instrução & categoria & 207.901 & 4,86 & 1,73 & 1,00 & 7,00 \\
Rendimento habitual & R\$ & 204.142 & 3.417,54 & 4.795,91 & 10,00 & 250.000,00 \\
Log rendimento habitual & log R\$ & 204.142 & 7,74 & 0,84 & 5,30 & 9,95 \\
Horas habituais & horas/sem. & 207.901 & 40,13 & 11,48 & 1,00 & 120,00 \\
Mulheres & \% & 207.901 & 43,73 & 49,61 & 0,00 & 100,00 \\
Brancos & \% & 207.901 & 43,17 & 49,53 & 0,00 & 100,00 \\
Trabalhadores formais & \% & 207.901 & 42,38 & 49,42 & 0,00 & 100,00 \\
\bottomrule
\end{tabular}
}
\end{center}

\vspace{0.15cm}

\textbf{Leitura:} esta tabela resume o corte transversal usado para
descrever a exposição à IA. As médias são ponderadas pelo peso da PNADc,
portanto representam a força de trabalho ocupada, não apenas a amostra
observada.
\end{frame}

\begin{frame}{Evidência Visual Preliminar --- Distribuição}
\phantomsection\label{eviduxeancia-visual-preliminar-distribuiuxe7uxe3o}
\begin{columns}[T]
\begin{column}{0.58\linewidth}
\includegraphics[width=1\linewidth,height=\textheight,keepaspectratio]{figures/fig_01_histograma_kde_exposicao.png}
\end{column}

\begin{column}{0.42\linewidth}
\textbf{O que o gráfico mostra?}

\begin{itemize}
\tightlist
\item
  A distribuição é assimétrica: muita massa em baixa exposição.
\item
  Há uma cauda relevante de ocupações com exposição alta.
\item
  A média sozinha esconde a heterogeneidade entre tipos de ocupação.
\end{itemize}
\end{column}
\end{columns}
\end{frame}

\begin{frame}{Evidência Visual Preliminar --- Gradientes}
\phantomsection\label{eviduxeancia-visual-preliminar-gradientes}
\begin{columns}[T]
\begin{column}{0.58\linewidth}
\includegraphics[width=1\linewidth,height=\textheight,keepaspectratio]{figures/fig_01_populacao_por_gradiente_ilo.png}
\end{column}

\begin{column}{0.42\linewidth}
\textbf{Como interpretar os gradientes?}

\begin{itemize}
\tightlist
\item
  Gradientes 1 e 2 indicam exposição parcial, mais próxima de
  complementaridade.
\item
  Gradientes 3 e 4 indicam exposição alta e mais consistente entre
  tarefas.
\item
  O grupo G3-G4 é a base conceitual para definir ocupações de alta
  exposição nas etapas causais.
\end{itemize}
\end{column}
\end{columns}
\end{frame}

\begin{frame}{Evidência Visual Preliminar --- Perfil Demográfico}
\phantomsection\label{eviduxeancia-visual-preliminar-perfil-demogruxe1fico}
\begin{columns}[T]
\begin{column}{0.6\linewidth}
\includegraphics[width=1\linewidth,height=\textheight,keepaspectratio]{figures/fig_04_demograficos_exposicao_3grupos.png}
\end{column}

\begin{column}{0.4\linewidth}
\textbf{Padrão substantivo}

\begin{itemize}
\tightlist
\item
  Mulheres, brancos, jovens e mais escolarizados aparecem mais em alta
  exposição.
\item
  Esse padrão não é individual; reflete principalmente sorting
  ocupacional.
\item
  A etapa 2 parte exatamente desse grupo de ocupações mais expostas.
\end{itemize}
\end{column}
\end{columns}
\end{frame}

\section{Etapa 2: CAGED + ILO}\label{etapa-2-caged-ilo}

\begin{frame}{Por que CAGED + Diferenças em Diferenças?}
\phantomsection\label{por-que-caged-diferenuxe7as-em-diferenuxe7as}
A etapa 1 é descritiva: mostra \textbf{quem está exposto}, mas não
identifica efeito causal.

A etapa 2 muda a pergunta:

\begin{quote}
Após o lançamento do ChatGPT, ocupações altamente expostas à IA tiveram
mudanças diferenciais em salário ou volume de admissões?
\end{quote}

\begin{itemize}
\tightlist
\item
  O CAGED mede fluxos mensais do mercado formal.
\item
  O evento é o lançamento do ChatGPT em novembro de 2022.
\item
  O tratamento vem da exposição ocupacional medida pelo índice ILO.
\end{itemize}
\end{frame}

\begin{frame}{Etapa 2 --- Ficha Técnica do Painel}
\phantomsection\label{etapa-2-ficha-tuxe9cnica-do-painel}
\begin{center}
\resizebox{0.78\textwidth}{!}{\scriptsize
\begin{tabular}{lr}
\toprule
Item & Valor \\
\midrule
Fonte & Novo CAGED + ILO WP140 \\
Janela analítica & 2021-01 a 2025-06 \\
Movimentações CAGED agregadas & 198,8 milhões \\
Células ocupação-mês do painel & 32.988 \\
Ocupações CBO 4d & 629 \\
Períodos mensais & 54 \\
Ocupações tratadas & 131 \\
Participação tratada & 20,8\% \\
Evento & Lançamento do ChatGPT (nov/2022) \\
\bottomrule
\end{tabular}
}
\end{center}

\vspace{0.2cm}

\begin{itemize}
\tightlist
\item
  Unidade da regressão: ocupação CBO 4 dígitos × mês.
\item
  As 32.988 células ocupação-mês resumem 198,8 milhões de movimentações
  individuais do CAGED.
\item
  Tratamento de extremos: salários de admissão limitados aos percentis 1
  e 99; logs recalculados após o ajuste.
\item
  Grupo tratado: ocupações no top 20\% do score ILO.
\item
  Grupo controle: demais ocupações.
\end{itemize}
\end{frame}

\begin{frame}{Etapa 2 --- Estatísticas Descritivas dos Outcomes}
\phantomsection\label{etapa-2-estatuxedsticas-descritivas-dos-outcomes}
\begin{center}
\resizebox{0.96\textwidth}{!}{\scriptsize
\begin{tabular}{lrrrrrr}
\toprule
Variável & Unidade & N & Média & Desv. Pad. & Min & Max \\
\midrule
Admissões & vagas/mês & 32.988 & 3.153,53 & 13.879,98 & 0,00 & 289.900,00 \\
Desligamentos & vagas/mês & 32.988 & 2.873,65 & 12.450,54 & 0,00 & 284.338,00 \\
Saldo líquido & vagas/mês & 32.988 & 279,88 & 2.323,00 & -46.650,00 & 90.556,00 \\
Salário admissão & R\$ & 32.988 & 6.131,12 & 158.475,43 & 0,00 & 18.799.538,19 \\
Salário admissão & SM & 32.988 & 4,74 & 120,53 & 0,00 & 15.511,17 \\
Idade média & anos & 32.988 & 33,23 & 5,43 & 0,00 & 73,00 \\
Mulheres nas admissões & \% & 32.988 & 0,00 & 0,00 & 0,00 & 0,00 \\
Superior nas admissões & \% & 32.988 & 21,47 & 26,43 & 0,00 & 100,00 \\
Score ILO & 0-1 & 32.988 & 0,28 & 0,12 & 0,12 & 0,63 \\
\bottomrule
\end{tabular}
}
\end{center}

\vspace{0.15cm}

\textbf{Leitura:} as variáveis de quantidade medem fluxos de emprego; as
variáveis salariais medem preço do trabalho no momento da admissão, a
margem mais limpa para capturar mudanças recentes de demanda.
\end{frame}

\begin{frame}{Estratégia Empírica --- Modelo da Etapa 2}
\phantomsection\label{estratuxe9gia-empuxedrica-modelo-da-etapa-2}
\textbf{Equação de regressão}

\[
y_{ot} =
\beta\,(\text{Pós}_t \times \text{AltaExp}_o)
+ \gamma X_{ot}
+ \alpha_o
+ \delta_t
+ \varepsilon_{ot}
\]

\vspace{0.3cm}

\begin{itemize}
\tightlist
\item
  \(o\) indexa ocupações CBO 4 dígitos.
\item
  \(t\) indexa meses.
\item
  A variação de interesse combina \textbf{antes/depois} do ChatGPT com
  \textbf{alta/baixa exposição} ocupacional.
\end{itemize}
\end{frame}

\begin{frame}{Etapa 2 --- Legenda Detalhada do Modelo}
\phantomsection\label{etapa-2-legenda-detalhada-do-modelo}
\begin{itemize}
\tightlist
\item
  \(y_{ot}\) é o outcome da ocupação \(o\) no mês \(t\): log do salário
  real, log das admissões, desligamentos ou composição das admissões.
\item
  \(\text{AltaExp}_o\) vale 1 para ocupações no top 20\% do score ILO.
  Ela captura a diferença \textbf{transversal} entre ocupações mais e
  menos expostas.
\item
  \(\text{Pós}_t\) vale 1 a partir de dezembro de 2022. Ela captura a
  diferença \textbf{temporal} após o lançamento do ChatGPT.
\item
  \(\beta\) é o coeficiente de interesse: o efeito diferencial
  pós-ChatGPT nas ocupações altamente expostas, sob tendências
  paralelas.
\item
  \(X_{ot}\) inclui controles que mudam ao longo do tempo, como idade
  média, participação feminina e participação com superior.
\end{itemize}
\end{frame}

\begin{frame}{Etapa 2 --- Efeitos Fixos e Erros-Padrão}
\phantomsection\label{etapa-2-efeitos-fixos-e-erros-padruxe3o}
\begin{itemize}
\tightlist
\item
  \(\alpha_o\) é o efeito fixo de ocupação. Ele absorve diferenças
  permanentes entre ocupações, como prestígio, complexidade técnica e
  nível salarial estrutural.
\item
  \(\delta_t\) é o efeito fixo de mês. Ele absorve choques comuns a todo
  o mercado formal, como inflação, sazonalidade e ciclo macroeconômico.
\item
  \(\varepsilon_{ot}\) é o erro idiossincrático.
\item
  Os erros-padrão são clusterizados por CBO 4d, porque o tratamento
  varia no nível da ocupação.
\item
  A hipótese central é: sem o ChatGPT, ocupações tratadas e controle
  seguiriam trajetórias paralelas.
\end{itemize}
\end{frame}

\begin{frame}{Etapa 2 --- Balanço Pré-Tratamento}
\phantomsection\label{etapa-2-balanuxe7o-pruxe9-tratamento}
\begin{center}
\resizebox{0.90\textwidth}{!}{\scriptsize
\begin{tabular}{lrrrr}
\toprule
Variável & Controle & Tratamento & Dif. norm. & Balanceado \\
\midrule
Admissões & 2.655,34 & 3.409,84 & 0,06 & Sim \\
Desligamentos & 2.321,33 & 2.916,92 & 0,06 & Sim \\
Saldo líquido & 334,01 & 492,92 & 0,09 & Sim \\
Salário médio (R\$) & 4.073,21 & 5.919,86 & 0,19 & Sim \\
Idade média & 32,75 & 31,80 & -0,17 & Sim \\
\% mulheres & 0,00 & 0,00 & -- & Não \\
\% superior & 15,57 & 41,89 & 1,15 & Não \\
\bottomrule
\end{tabular}
}
\end{center}

\vspace{0.15cm}

\textbf{Leitura:} as variáveis de fluxo são mais comparáveis no
pré-tratamento. Salário, escolaridade e composição feminina são
desbalanceados porque a própria exposição à IA está concentrada em
ocupações \emph{white collar}. Os efeitos fixos de ocupação tratam
diferenças permanentes.
\end{frame}

\section{Etapa 3: CAGED + Anatel}\label{etapa-3-caged-anatel}

\begin{frame}{Por que Adicionar Conectividade Municipal?}
\phantomsection\label{por-que-adicionar-conectividade-municipal}
A IA generativa é uma tecnologia digital. Mesmo que duas ocupações
tenham a mesma exposição técnica, a difusão pode ser maior onde há
melhor infraestrutura de internet.

A etapa 3 testa se o canal ocupacional da etapa 2 é amplificado em
municípios com maior conectividade.

\begin{itemize}
\tightlist
\item
  Unidade: município × ocupação × mês.
\item
  Fonte de conectividade: indicadores municipais da Anatel.
\item
  Estratégia: diferença-em-diferenças-em-diferenças.
\end{itemize}
\end{frame}

\begin{frame}{Etapa 3 --- Ficha Técnica do Painel Municipal}
\phantomsection\label{etapa-3-ficha-tuxe9cnica-do-painel-municipal}
\begin{center}
\resizebox{0.82\textwidth}{!}{\scriptsize
\begin{tabular}{lr}
\toprule
Item & Valor \\
\midrule
Fonte & Novo CAGED + Anatel + IBGE + ILO WP140 \\
Unidade do painel & Município × ocupação × mês \\
Janela analítica & 2021-01 a 2025-06 \\
Movimentações CAGED agregadas & 161,9 milhões \\
Células município-ocupação-mês & 5.534.808 \\
Municípios & 657 \\
Ocupações CBO 4d & 614 \\
UFs & 28 \\
Alta conectividade & 80,5\% \\
\bottomrule
\end{tabular}
}
\end{center}

\vspace{0.2cm}

\begin{itemize}
\tightlist
\item
  O painel municipal resume 161,9 milhões de movimentações em 5.534.808
  células município-ocupação-mês.
\item
  Essa granularidade permite comparar ocupações semelhantes em ambientes
  digitais diferentes.
\item
  Tratamento de extremos: salário nominal e salário real de admissão
  limitados aos percentis 1 e 99.
\item
  Alta conectividade é definida antes do tratamento, evitando que o
  próprio choque pós-2022 redefina o grupo.
\end{itemize}
\end{frame}

\begin{frame}{Etapa 3 --- Estatísticas Descritivas}
\phantomsection\label{etapa-3-estatuxedsticas-descritivas}
\begin{center}
\resizebox{0.96\textwidth}{!}{\scriptsize
\begin{tabular}{lrrrrrr}
\toprule
Variável & Unidade & N & Média & Desv. Pad. & Min & Max \\
\midrule
Admissões & vagas/mês & 5.534.808 & 15,30 & 155,07 & 0,00 & 34.019,00 \\
Desligamentos & vagas/mês & 5.534.808 & 13,96 & 141,03 & 0,00 & 36.875,00 \\
Saldo líquido & vagas/mês & 5.534.808 & 1,35 & 33,95 & -7.322,00 & 11.635,00 \\
Salário real admissão & R\$ dez/2024 & 5.534.808 & 5.482,50 & 3.010.047,82 & 0,00 & 5.505.546.848,14 \\
Idade média & anos & 5.534.808 & 26,42 & 15,20 & 0,00 & 94,00 \\
Mulheres nas admissões & \% & 5.534.808 & 26,88 & 37,94 & 0,00 & 100,00 \\
Superior nas admissões & \% & 5.534.808 & 13,55 & 29,57 & 0,00 & 100,00 \\
Jovens nas admissões & \% & 5.534.808 & 34,63 & 37,20 & 0,00 & 100,00 \\
Score ILO & 0-1 & 5.534.808 & 0,28 & 0,14 & 0,12 & 0,63 \\
Penetração banda larga & \% & 5.534.808 & 64,34 & 24,61 & 2,06 & 127,27 \\
Fibra no pré & \% & 5.534.808 & 0,00 & 0,00 & 0,00 & 0,00 \\
PIB per capita & R\$ & 5.534.808 & 52.619,04 & 47.935,23 & 8.167,85 & 473.953,85 \\
\bottomrule
\end{tabular}
}
\end{center}
\end{frame}

\begin{frame}{Etapa 3 --- Conectividade e Amostra}
\phantomsection\label{etapa-3-conectividade-e-amostra}
\begin{center}
\resizebox{0.82\textwidth}{!}{\scriptsize
\begin{tabular}{lrrr}
\toprule
Variável & Baixa conect. & Alta conect. & Diferença \\
\midrule
Admissões & 5,39 & 16,27 & 10,87 \\
Salário real admissão & 5.580,02 & 5.473,01 & -107,01 \\
Idade média & 24,27 & 26,63 & 2,36 \\
\% mulheres & 21,87 & 27,37 & 5,50 \\
\% superior & 8,81 & 14,01 & 5,19 \\
Score ILO & 0,29 & 0,28 & -0,00 \\
Banda larga fixa & 17,80 & 68,87 & 51,07 \\
Fibra & 0,00 & 0,00 & 0,00 \\
\bottomrule
\end{tabular}
}
\end{center}

\vspace{0.15cm}

\textbf{Leitura:} municípios de alta conectividade são estruturalmente
diferentes. Por isso, a etapa 3 não compara apenas municípios conectados
contra não conectados; ela usa a interação com exposição ocupacional e
período pós-ChatGPT.
\end{frame}

\begin{frame}{Estratégia Empírica --- Modelo da Etapa 3}
\phantomsection\label{estratuxe9gia-empuxedrica-modelo-da-etapa-3}
\textbf{Equação Triple-DiD}

\[
\begin{aligned}
y_{mot} =\;&
\beta_1(\text{Pós}_t \times \text{AltaExp}_o)
+ \beta_2(\text{Pós}_t \times \text{AltaConect}_m) \\
&+ \beta_3(\text{AltaExp}_o \times \text{AltaConect}_m)
+ \beta_4(\text{Pós}_t \times \text{AltaExp}_o \times \text{AltaConect}_m) \\
&+ \gamma X_{mot} + \alpha_{mo} + \delta_t + \varepsilon_{mot}
\end{aligned}
\]

\vspace{0.2cm}

O foco é \(\beta_4\): o efeito adicional em ocupações expostas, no
pós-ChatGPT, dentro de municípios mais conectados.
\end{frame}

\begin{frame}{Etapa 3 --- Legenda Detalhada do Triple-DiD}
\phantomsection\label{etapa-3-legenda-detalhada-do-triple-did}
\begin{itemize}
\tightlist
\item
  \(m\) indexa municípios, \(o\) ocupações e \(t\) meses.
\item
  \(\text{AltaConect}_m\) vale 1 para municípios acima do ponto de corte
  de conectividade no pré-tratamento.
\item
  \(\beta_1\) captura o canal ocupacional médio da etapa 2.
\item
  \(\beta_2\) captura mudanças pós-ChatGPT em municípios conectados para
  todas as ocupações.
\item
  \(\beta_3\) captura diferenças estruturais entre ocupações expostas
  localizadas em municípios conectados.
\item
  \(\beta_4\) isola a amplificação digital: exposição ocupacional alta
  \textbf{mais} capacidade local de difusão da tecnologia.
\end{itemize}
\end{frame}

\begin{frame}{Etapa 3 --- Efeitos Fixos e Identificação}
\phantomsection\label{etapa-3-efeitos-fixos-e-identificauxe7uxe3o}
\begin{itemize}
\tightlist
\item
  \(\alpha_{mo}\) é o efeito fixo de município × ocupação.
\item
  Ele compara a mesma ocupação dentro do mesmo município ao longo do
  tempo.
\item
  Isso absorve diferenças permanentes, como estrutura produtiva local,
  nível salarial local e composição ocupacional fixa.
\item
  \(\delta_t\) remove choques agregados comuns a todos os municípios e
  ocupações.
\item
  A identificação vem da diferença relativa pós-ChatGPT entre ocupações
  expostas e não expostas, em municípios com alta e baixa conectividade.
\end{itemize}
\end{frame}

\begin{frame}{Encerramento}
\phantomsection\label{encerramento}
\textbf{Mensagem central da apresentação}

\begin{itemize}
\tightlist
\item
  A etapa 1 mostra a geografia social da exposição à IA no Brasil.
\item
  A etapa 2 transforma essa exposição em tratamento econométrico.
\item
  A etapa 3 testa o mecanismo de difusão: a infraestrutura digital pode
  amplificar o choque.
\item
  As tabelas descritivas e de amostra são o elo entre dados brutos,
  especificação e interpretação.
\end{itemize}
\end{frame}




\end{document}
